\documentclass[10pt, a4paper]{article}

\usepackage{homework}

\title{\textbf{《数值代数》上机作业15}}
\author{岳瀚阳\quad\quad\quad\quad 3230104443}
\date{\today}

\begin{document}

\maketitle

%----------------------------------------------------------------------
\section{经典 Jacobi 方法}
%----------------------------------------------------------------------

Jacobi 方法是求实对称矩阵全部特征值和特征向量的经典算法。
其核心思想是：通过一系列平面上（Givens）旋转，逐次对消非对角元，最终将矩阵对角化。
设 $A^{(0)} = A$，第 $k$ 步选取非对角元 $a_{pq}^{(k)}$，构造 Jacobi 旋转矩阵
\[
J(p,q,\theta) = \begin{bmatrix}
1 & & & & & \\
& \ddots & & & & \\
& & c & \cdots & -s & \\
& & \vdots & \ddots & \vdots & \\
& & s & \cdots & c & \\
& & & & & \ddots & \\
& & & & & & 1
\end{bmatrix},
\qquad c = \cos\theta,\; s = \sin\theta,
\]
其中 $J$ 仅在 $(p,p), (q,q)$ 处为 $c$，$(p,q)$ 处为 $-s$，$(q,p)$ 处为 $s$，其余位置同单位阵。
取 $A^{(k+1)} = J^{\!\top} A^{(k)} J$，则 $a_{pq}^{(k+1)} = 0$。

旋转角度 $\theta$ 满足
\[
\tan 2\theta = \frac{2a_{pq}}{a_{pp} - a_{qq}},
\]
为避免直接计算反三角函数，采用
\[
\tau = \frac{a_{qq} - a_{pp}}{2a_{pq}}, \qquad
t = \frac{\mathrm{sgn}(\tau)}{|\tau| + \sqrt{1 + \tau^2}}, \qquad
c = \frac{1}{\sqrt{1 + t^2}}, \quad s = c t.
\]

实际计算时，不必每步寻找全局最大非对角元。采用逐轮扫描策略：每轮遍历所有上三角位置，仅对满足 $|a_{pq}| > \eta$ 的元素施以旋转。
初始阈值取为
\[
\eta^{(0)} = \frac{1}{n}\sqrt{\sum_{i<j} |a_{ij}|^2},
\]
每轮扫描结束后重估 $\eta$，当 $\eta < \varepsilon$ 时停止。
特征向量由所有旋转的累积给出：$V^{(k+1)} = V^{(k)} J$。

%----------------------------------------------------------------------
\section{数值结果}
%----------------------------------------------------------------------

用经典 Jacobi 方法计算教材第 244 页 1(2) 的矩阵：
\[
A = \begin{bmatrix}
4 & 1 & & & \\
1 & 4 & 1 & & \\
& \ddots & \ddots & \ddots & \\
& & 1 & 4 & 1 \\
& & & 1 & 4
\end{bmatrix}_{n\times n},
\qquad n = 50, 60, \ldots, 100.
\]

该矩阵的特征值有显式表达式 $\lambda_k = 4 + 2\cos\frac{k\pi}{n+1}$，可用于检验算法精度。

结果汇总见表~\ref{tab:jacobi_results}。

\begin{table}[H]
    \centering
    \caption{经典 Jacobi 方法的数值结果}
    \label{tab:jacobi_results}
    \begin{tabular}{cccccc}
        \toprule
        $n$ & Sweeps & $\lambda_{\min}$ & $\lambda_{\max}$ &
        特征值最大差异 & 特征向量最大残差 \\
        \midrule
        50  & 16 & 2.0037933425 & 5.9962066575 & $4.53\times10^{-14}$ & $2.82\times10^{-8}$ \\
        60  & 200 & 2.0026518202 & 5.9973481798 & $6.04\times10^{-14}$ & $1.62\times10^{-8}$ \\
        70  & 16 & 2.0019575470 & 5.9980424530 & $7.28\times10^{-14}$ & $4.51\times10^{-8}$ \\
        80  & 16 & 2.0015040950 & 5.9984959050 & $1.07\times10^{-13}$ & $7.49\times10^{-8}$ \\
        90  & 17 & 2.0011917189 & 5.9988082811 & $1.08\times10^{-13}$ & $1.34\times10^{-8}$ \\
        100 & 18 & 2.0009674354 & 5.9990325646 & $9.15\times10^{-14}$ & $1.21\times10^{-8}$ \\
        \bottomrule
    \end{tabular}
\end{table}

所有矩阵的特征值与 numpy.eigvalsh 的最大差异均在 $10^{-13}$ 量级，达到双精度下 Jacobi 方法的典型精度。
最小与最大特征值与理论值完全一致（小数点后 10 位全同）。

以 $n = 50$ 为例，特征向量正交性误差 $\|V^{\!\top} V - I\|_2 \approx 3.72\times 10^{-14}$。
最小特征值对应的特征向量前 5 个分量近似交替符号的正弦波：
\[
v_{\min} \approx [0.01219,\; -0.02434,\; 0.03639,\; -0.04830,\; 0.06003]^{\!\top},
\]
最大特征值对应的特征向量则为单调正弦波：
\[
v_{\max} \approx [0.01219,\; 0.02434,\; 0.03639,\; 0.04830,\; 0.06003]^{\!\top},
\]
这与理论特征向量 $v_k(j) = \sin(jk\pi/(n+1))$ 完全吻合。

%----------------------------------------------------------------------
\section{结果分析}
%----------------------------------------------------------------------

\textbf{（1）收敛性与扫描轮数。}
逐轮扫描策略每轮遍历所有 $O(n^2)$ 个非对角元，对大于阈值的元素施以 Jacobi 旋转。
$n = 50$ 时需 16 轮，$n = 100$ 时需 18 轮，轮数与矩阵规模的相关性较弱。
但 $n = 60$ 时达到了上限 200 轮。这是因为经过十余轮扫描后，该矩阵的非对角元范数恰好降至约 $10^{-12}$ 量级，离容限 $10^{-14}$ 差两个数量级，却又不足以被 Jacobi 旋转有效削减，导致算法反复扫描直至触及上限。此现象与特定矩阵规模下舍入误差的偶然叠加有关，不影响特征值的最终精度。

\textbf{（2）特征值精度与特征向量精度。}
Jacobn 方法的特征值收敛具有二次（甚至三次）速度，特征值精度可达机器精度量级。
然而特征向量的精度受到累积舍入误差的影响：每次旋转需更新整行和整列，$n$ 轮扫描共涉及 $O(n^3)$ 次浮点运算，舍入误差累积导致特征向量残差量级约为 $10^{-8}$。
通过重新正交化可进一步提高特征向量精度，但就本题而言，已足以满足工程需求。

\textbf{（3）与隐式对称 QR 算法的比较。}
在作业 14 中，隐式对称 QR 算法对同一矩阵特征向量的残差量级为 $10^{-12}$，优于 Jacobi 方法约 4 个数量级。
这是因为 Householder 三对角化保留了矩阵的带状结构，而 Jacobi 方法在旋转过程中引入大量 fill-in，矩阵从三对角迅速变为稠密，增加了运算量和误差累积。
然而，Jacobi 方法实现简明，适合小规模矩阵的教学演示。

%----------------------------------------------------------------------
\appendix
\section{源代码}
%----------------------------------------------------------------------

\subsection{\texttt{hw15\_test.py}（主程序）}
\lstinputlisting{../tests/hw15_test.py}

\subsection{\texttt{jacobi\_eigen.py}（经典 Jacobi 方法）}
\lstinputlisting{../algorithms/jacobi_eigen.py}

\end{document}
